% !TeX program = xelatex
% ============================================================
% 符号说明（问题一独立文件，按 CUMCM-LaTeX-Template 组织）
%
% 本文件只列出问题一已经使用的符号。
% 后续完成问题二、问题三后，在 \bottomrule 之前增加对应符号，
% 再将 \section{符号说明} 及其表格并入总论文。
% ============================================================

\PassOptionsToPackage{quiet}{xeCJK}
\documentclass[withoutpreface,notoc]{cumcmthesis}

\RenewDocumentCommand{\textbf}{m}{{\heiti #1}}

\usepackage{amsmath,amssymb,bm}
\usepackage{array}
\usepackage{tabularx}
\usepackage{booktabs}
\usepackage{etoolbox}
\usepackage{enumitem}

\BeforeBeginEnvironment{tabular}{\zihao{-5}}
\BeforeBeginEnvironment{tabularx}{\zihao{-5}}

\newcolumntype{L}{>{\raggedright\arraybackslash}X}
\newcolumntype{C}{>{\centering\arraybackslash}X}

\newcommand{\bs}{\boldsymbol{S}}
\newcommand{\bp}{\boldsymbol{p}}
\newcommand{\bu}{\boldsymbol{u}}
\newcommand{\ba}{\boldsymbol{a}}
\newcommand{\bA}{\boldsymbol{A}}
\newcommand{\bB}{\boldsymbol{B}}
\newcommand{\bC}{\boldsymbol{C}}
\newcommand{\bO}{\boldsymbol{O}}
\newcommand{\Ball}{\mathcal{B}}
\newcommand{\T}{\mathrm{T}}
\newcommand{\R}{\mathbb{R}}
\DeclareMathOperator{\diam}{diam}
\DeclareMathOperator{\conv}{conv}

\title{符号说明}

\begin{document}
\maketitle

\section{符号说明}

为方便阅读，本文使用的问题一主要符号及其说明如\cref{tab:符号说明}所示。
跨问题复用的符号只在本表中定义一次；问题二、问题三的专用变量待模型确定后统一补充。

\begin{table}[H]
  \centering
  \caption{问题一主要符号及其含义}
  \label{tab:符号说明}
  \begin{tabularx}{\textwidth}{>{\centering\arraybackslash}p{0.17\textwidth}
                                L
                                >{\centering\arraybackslash}p{0.18\textwidth}}
    \toprule
    符号 & 含义 & 单位或类型 \\
    \midrule
    $n$ & 同一干扰源的有效示向度观测数 & 个 \\
    $\bs_i=(x_i,y_i)$ & 第 $i$ 个检测点的平面坐标 & m \\
    $\theta_i$ & 第 $i$ 个检测点测得的示向度 & 度 \\
    $\alpha$ & 测向误差半角，本题取 $\alpha=1^\circ$ & 度 \\
    $\bp=(x,y)$ & 干扰源候选位置 & m \\
    $\bu_i^-,\bu_i^+$ & 第 $i$ 个前向楔形的下、上边界单位方向向量 & 无量纲 \\
    $W_i$ & 第 $i$ 个观测形成的闭前向楔形 & 点集 \\
    $\ba_k,b_k$ & 第 $k$ 条半平面约束的系数向量与右端常数 & 无量纲 \\
    $P$ & 全部前向楔形的交会定位区域 & 点集 \\
    $V(P)$ & $P$ 的逆时针凸包顶点序列 & 点集 \\
    $D$ & 有界定位区域的直径 & m \\
    $\bA,\bB$ & 一组最远点，满足 $\|\bA-\bB\|_2=D$ & 点 \\
    $\bC$ & 最远点对的中点，$\bC=(\bA+\bB)/2$ & 点 \\
    $\Ball(o,r)$ & 以 $o$ 为圆心、$r$ 为半径的闭圆盘 & 点集 \\
    $R_{\min}$ & 定位区域顶点集的最小覆盖圆半径 & m \\
    $\bO$ & 目标圆域的中心，取为坐标原点 & 点 \\
    $\delta_i$ & 检测点 $\bs_i$ 到定位区域 $P$ 的最短距离 & m \\
    $P_{\mathrm{phys}}$ & 加入目标圆域、接收范围和过近条件后的补充修正区域 & 点集 \\
    \bottomrule
  \end{tabularx}
\end{table}

% ============================================================
% 合并到总论文时：
%   1. 保留 \section{符号说明} 及 table 环境；
%   2. 在 \bottomrule 之前按问题分组增加新符号；
%   3. 同名符号只保留一次，并在含义中说明适用范围；
%   4. 不把本文件的 \maketitle 复制到总论文正文中。
% ============================================================

\end{document}
